% Modernised reading — fonds Alexandre Grothendieck, shelfmark 131
% (pages 1-41, the whole folder).
%
% This is an INTERPRETATION, not a transcription. It re-reads the pages in
% today's notation and vocabulary, states the hypotheses the manuscript leaves
% implicit, and drops the critical apparatus so that the mathematics reads
% straight through. Everything the brackets and underlines carried moves into a
% footnote. It is held to being mathematically correct as it stands; where the
% manuscript is loose or mistaken, the true statement is what appears here and
% the footnote says what the page has. In any dispute of substance the
% transcription governs: whoever cites Grothendieck must cite batch-01.fr.tex
% (pp. 1-20) or batch-02.fr.tex (pp. 21-40).
%
% Pass: Opus 5.5 (claude-opus-5-5), 2026-09-24 — first pass, unchecked against the pages by a human.
% The folder was read whole, its three batches together; this is one document
% for the shelfmark. It was made on Opus 5.5 and has not been measured against
% the reference reading of folder 115 (made on Opus 5). The literature named
% in the footnotes is cited from memory and was not checked against the
% sources.
%
% Scope. Batch 3 is page 41 alone, a leaf of a foreign typescript with nothing
% in his hand; its transcription has an empty body and no \page, and nothing
% here reads it (no \pagerange names page 41). The same holds for pages 7, 9,
% 12, 20 and the odd pages 23-39 (typed leaves by others, « n° 371 » and
% « n° 392 »), and for the blank pages 10, 18, 21. Page 17 carries two words.
%
% Spine. One theorem and its tools, written in the reverse of their logical
% order. Pages 22-40 are the « notes précédentes » that pages 3-5 cite: the
% « cor 2 à prop. 1 » of page 3 is the Corollaire 2 of page 28, the « prop. 2
% des notes » is the Proposition 2 of pages 30-32, the « prop. 3 des notes
% préc. » and its « condition de constructibilité » are the Prop. 3 of page
% 32. The identification is ours, from the numbering and the content; the
% pages do not say where the earlier notes are. The two marginal programmes of
% page 32 (« Transformer 4°) en une [condition] plus facile à vérifier, avec
% ... dim 1 » and « Transformer en une condition de modularité générique »)
% are what the « critère de modularité a) locale / b) générique » of page 1
% (conditions 7°, 8°) carry out.
%   pp. 1-2    first list of eight conditions
%   pp. 3-5    proof by induction on dim S, citing pp. 28, 30-32, 32
%   pp. 6, 8   axioms (A), (B), (C) (page 6 reformulates page 8's C_1)
%   p. 11      the list for the unramified case, with (9°) = (RG)
%   pp. 13-16  the Théorème, corollaries, (RG), examples, applications
%   p. 19      proof of (A) and (B)
%   pp. 22-40  Prop. 1, Cor. 1, Cor. 2, Prop. 2, Prop. 3, modular pairs,
%              proof of Prop. 3
%
% Notation fixed for the folder. « préschéma » is written « schéma ». The
% conditions of the Théorème are numbered (1)-(8) as on page 13; page 1's
% numbering is translated through a concordance given in the spine section,
% and page 1's 4° (no counterpart on page 13) is called (E). The conditions of
% pages 22-32 are written (L1)-(L4), (R_loc), (R_0): the arrows he draws
% under the L and the two small circles on L_3, L_4 are dropped (footnoted).
% « Modulaire » is the notion of page 34, fixed in the spine section and used
% from page 1 on.
%
% Substantive departures, each footnoted where it occurs:
%   - p. 15, Cor. 2 b): the bracket « tel que I_p = 0 » makes the condition
%     fail already for a closed immersion (counterexample given); the reading
%     states it with Spec A_p schematically dense (page 6's marginal
%     « d'adhérence schématique ») and I only nilpotent;
%   - p. 34: « Hom_k(k(t'), k(t)) » has its variance reversed; the reading
%     writes Hom_k(k(t), k(t'));
%   - p. 26: the step u'p_1 = u'p_2 invokes the injectivity of 1°) for
%     T' x_T T', which is not local essentially of finite type; the reading
%     supplies the extension of injectivity through (L2);
%   - p. 30: the construction of xi_1^* needs F(B^) -> lim F(B_n) surjective,
%     which (L4) as stated on p. 24 (injective) does not give; the reading
%     assumes bijectivity, as condition (4) of page 13 does;
%   - p. 4, step (3): effectivity of the descent datum is asserted; the
%     reading says what makes it hold in the quasi-compact case and flags the
%     rest;
%   - p. 16: « un schéma abélien ... est projectif » is false in general
%     (Raynaud); the reading says so;
%   - p. 38: the Lemme's « eta in F(S) » is read eta in F(Z);
%   - pp. 8 and 19: fibre products over S of S'-schemes are read over S'.
\documentclass[11pt,a4paper]{article}
% Le préambule commun à toutes les transcriptions du fonds Grothendieck.
%
% Il tient en une idée : **l'incertitude se compose, elle ne se résout pas.**
% Une transcription qui lisse un mot douteux a détruit la seule chose pour
% laquelle on la faisait — un lecteur doit pouvoir savoir, à chaque endroit,
% ce qui était sur le feuillet et ce qui a été ajouté. Les six macros qui
% suivent sont donc l'appareil critique, et non de la décoration.
%
% Les mêmes commandes sont reconnues par `scripts/render.mjs`, qui produit la
% vue de lecture du site. Une macro ajoutée ici doit l'être aussi là-bas,
% faute de quoi le rendu échouera bruyamment — ce qui est voulu : un rendu
% silencieusement faux est pire qu'un rendu absent.

\ProvidesPackage{grothendieck}[2026/08/08 Transcriptions du fonds Grothendieck]

\usepackage{amsmath,amssymb,amsthm}
\usepackage{mathrsfs}
\usepackage{stmaryrd}
% Les diagrammes commutatifs sont partout dans ce fonds : c'est la langue
% même de la géométrie algébrique de Grothendieck, et une transcription qui
% les rendrait en prose ne transcrirait rien.
\usepackage{tikz-cd}
% Français par défaut — c'est la langue des feuillets — mais l'anglais est
% chargé aussi : la traduction et le résumé sont des documents anglais, et la
% typographie française (espaces avant les deux-points) y serait une faute.
% Ces documents basculent par \selectlanguage{english} en tête de corps.
\usepackage[english,french]{babel}
\usepackage{fontspec}
\usepackage{geometry}
\geometry{a4paper,margin=25mm,marginparwidth=28mm}
\usepackage{marginnote}
\usepackage{xcolor}
% ulem, not soul. `soul`'s \st and \ul are built on a character-by-character
% scan that XeLaTeX's font machinery breaks: the document fails with an
% undefined control sequence rather than degrading. ulem does the same two jobs
% and is Unicode-engine safe. `normalem` stops it hijacking \emph, which the
% transcriptions use for Grothendieck's own emphasis.
\usepackage[normalem]{ulem}
\usepackage{hyperref}
\hypersetup{colorlinks=true,linkcolor=black,urlcolor=black,pdfborder={0 0 0}}

% --- Métadonnées du lot -----------------------------------------------------
% Reprises telles quelles par le script de rendu, qui les affiche en tête de
% la vue de lecture. Elles fixent ce que le fichier prétend couvrir : sans
% elles, une transcription flotte sans point d'attache dans un fonds de seize
% mille feuillets.

\newcommand{\folder}[1]{\gdef\grothFolder{#1}}
\newcommand{\batch}[1]{\gdef\grothBatch{#1}}
\newcommand{\leaves}[2]{\gdef\grothFirst{#1}\gdef\grothLast{#2}}
\newcommand{\foldertitle}[1]{\gdef\grothFoldertitle{#1}}
\gdef\grothFolder{?}\gdef\grothBatch{?}\gdef\grothFirst{?}\gdef\grothLast{?}\gdef\grothFoldertitle{}

% --- Le résumé --------------------------------------------------------------
%
% Il ouvre la lecture modernisée, sous le titre « Résumé », et il est composé
% en petit. Ce n'est pas un ornement : c'est le seul passage du document qui
% ne soit pas des mathématiques, et un lecteur doit pouvoir le reconnaître
% avant de l'avoir lu — puis le sauter s'il connaît déjà le sujet, ou s'y
% arrêter s'il ne le connaît pas. Le filet qui le suit marque la reprise du
% corps.

\newenvironment{resume}
  {\small\setlength{\parskip}{0.45em}}
  {\par\medskip\hrule\medskip}

% --- Mots-clés ---------------------------------------------------------------
%
% En anglais, en clôture du résumé de la lecture modernisée : le vocabulaire
% moderne sous lequel chercher ce que les feuillets construisent. Le manifeste
% du site (`npm run manifest`) les extrait pour étiqueter la cote entière —
% cette ligne est la source unique des tags, il n'y a pas de fichier à côté.
\newcommand{\keywords}[1]{\par\smallskip\noindent
  {\footnotesize\textsc{Keywords} --- \textit{#1}}\par}

% --- L'appareil critique ----------------------------------------------------

% Le numéro de feuillet, en marge. C'est l'ancre qui permet de régler un
% désaccord sur une formule en regardant *un* feuillet plutôt qu'une chemise
% de six cents. Le site s'en sert aussi pour tourner les pages du fac-similé
% quand on fait défiler la transcription.
\newcommand{\leaf}[1]{%
  \marginnote{\footnotesize\color{gray!70}\textsf{#1}}%
  \label{leaf:#1}\ignorespaces}

% Illisible : on ne devine pas. Un mot inventé qui se lit comme les autres est
% le pire résultat possible.
\newcommand{\ill}{\textcolor{red!60!black}{[\,\ldots\,]}}

% Lecture douteuse : le mot est donné, et signalé comme incertain.
\newcommand{\uncertain}[1]{\uline{#1}}

% Ajout de l'éditeur — une lettre manquante, un mot sous-entendu.
\newcommand{\add}[1]{\textcolor{blue!50!black}{[#1]}}

% Biffé par Grothendieck lui-même. Ce qu'il a rayé fait partie du document :
% une rature dit souvent où la pensée a changé de direction.
\newcommand{\struck}[1]{\sout{#1}}

% Note du transcripteur, sur l'état matériel du feuillet.
\newcommand{\note}[1]{\par\smallskip{\small\color{gray!60!black}\textit{[#1]}}\par\smallskip}

% Note marginale de Grothendieck, distincte du corps du texte.
\newcommand{\marginal}[1]{\par\smallskip\noindent\hspace{1em}%
  {\small\color{gray!60!black}#1}\par\smallskip}

% --- Titre ------------------------------------------------------------------

\AtBeginDocument{%
  \begin{center}
    {\large\bfseries Fonds Alexandre Grothendieck --- cote n\textsuperscript{o}~\grothFolder}\\[2pt]
    {\itshape\grothFoldertitle}\\[4pt]
    {\small feuillets \grothFirst--\grothLast\ (lot \grothBatch)}\\[2pt]
    {\footnotesize\color{gray!60!black}Transcription établie d'après le fac-similé numérisé
     par l'Université de Montpellier.\\
     Les lectures incertaines sont soulignées, les illisibles notées [\,\ldots\,],
     les ajouts entre crochets.}
  \end{center}
  \vspace{1em}}

\endinput


\folder{131}
\pages{1}{41}
\dating{[vers 1971]}
\watermark{Édition de démonstration\\interprétation personnelle de l'œuvre}
\foldertitle{Foncteurs représentables par schéma quasi fini : notes manuscrites (s.d.) — lecture modernisée du dossier entier}

\begin{document}

\section*{Résumé}

\begin{resume}
Ces feuillets cherchent à reconnaître un espace géométrique à la seule manière
dont il répond quand on l'interroge.

Voici le problème. En géométrie algébrique, beaucoup d'objets qu'on voudrait
étudier — les courbes d'un certain type, les fibrés en droites sur une
variété, les morphismes entre deux variétés abéliennes — ne sont d'abord
connus que par un « questionnaire » : pour chaque espace test $T$, on sait dire
quelles sont les familles de ces objets paramétrées par $T$. Un tel
questionnaire s'appelle un foncteur $F$. On voudrait qu'il existe un vrai
espace $X$, un schéma, tel que se donner une famille paramétrée par $T$
revienne exactement à se donner une application de $T$ dans $X$ : on dit alors
que $X$ \emph{représente} $F$, et $X$ est l'« espace de modules » du problème.
Mais on ne peut presque jamais construire $X$ d'un coup. La question est donc :
quelles propriétés, vérifiables sur le questionnaire lui-même, garantissent
qu'un tel $X$ existe ?

Le dossier y répond dans un cas précis : celui où $X$ serait « discret
au-dessus de la base », c'est-à-dire n'aurait, au-dessus de chaque point, que
des points isolés en nombre fini localement — un schéma \emph{localement quasi
fini}. C'est le cas des foncteurs qu'il appelle non ramifiés : les morphismes
d'une variété abélienne dans une autre, les sous-groupes d'un groupe, les
composantes connexes des fibres d'un morphisme.

L'idée se comprend par une analogie avec les fonctions analytiques. Pour
savoir qu'une fonction est analytique, on regarde d'abord, en chaque point,
son développement de Taylor — une donnée purement formelle, qui ne dit rien de
ce qui se passe à distance finie ; puis on vérifie que la série converge sur
un petit disque ; puis on recolle les disques. Le dossier procède de même, en
trois étages. On interroge d'abord $F$ sur les « points épaissis »
(les anneaux artiniens) : si $F$ vient d'un espace $X$, les réponses doivent
reconstituer, point par point, un développement formel de $X$. On demande
ensuite que ce développement formel « converge » : qu'un paramétrage qui est
universel en un point le reste dans tout un voisinage. C'est ce que le dossier
appelle être \emph{modulaire}, et les deux conditions centrales du théorème —
un critère de modularité local et un critère générique — disent exactement que
la modularité se propage. Enfin des propositions « formelles » recollent les
morceaux. La démonstration monte alors en dimension : sur une base de
dimension zéro tout est formel ; sur une base locale complète on recolle le
point fermé au reste, déjà connu par récurrence ; on redescend du complété à
l'anneau local, puis de l'anneau local au voisinage.

Le dossier est un atelier, rangé à l'envers. Les outils — cinq énoncés de
recollement et une définition de la modularité — sont aux pages 22 à 40 ; le
théorème qui s'en sert est aux pages 1 à 19, et il les cite comme « les notes
précédentes ». En marge des outils, deux programmes sont notés : remplacer une
hypothèse difficile à vérifier par une condition « de dimension 1 », et une
hypothèse de constructibilité par une condition « générique ». Le théorème des
premières pages est l'exécution de ces deux programmes. Et à peine le
théorème écrit, les questions viennent en marge : lesquelles de ces conditions
sont vraiment nécessaires ? donner l'exemple type qui montre que l'une ne
découle pas des autres ; limiter la classe des anneaux sur lesquels il faut
les vérifier. C'est la manière d'un mathématicien qui ne tient pas un énoncé
pour acquis tant qu'il n'en a pas éprouvé chaque hypothèse.

Tout n'est pas établi : la dernière étape de la démonstration s'arrête sur
« on est ramené à vérifier la condition de constructibilité », et le
recollement final des pages 38 à 40 renvoie à « une autre proposition
formelle » qui n'est pas écrite. Les applications annoncées — la projectivité
des schémas abéliens, l'existence du schéma de Picard — ne sont qu'annoncées,
et la première, telle qu'elle est écrite, est fausse en général.

Pour aller plus loin, les noms d'aujourd'hui : critères de représentabilité
d'Artin, foncteur localement de présentation finie, pro-représentabilité et
critère de Schlessinger, effectivité des déformations formelles, espaces
algébriques et relations d'équivalence étales, foncteurs non ramifiés, schéma
de Picard.

\keywords{representable functor, locally quasi-finite morphism, Artin's
representability criteria, limit-preserving functor, pro-representable
functor, formal effectivity, fpqc descent, valuative criterion of
separatedness, unramified functor, étale morphism, open immersion,
constructible set, étale equivalence relation, algebraic space, Hom scheme
of abelian schemes, Picard scheme}
\end{resume}

\pagerange{1}{41}
\section*{Le fil du dossier, et les conventions}

\subsection*{Les stations}

Le dossier compte quarante et une pages, dont vingt-quatre de sa main ; les
autres sont blanches ou sont des feuillets de tapuscrits étrangers au sujet
(des rédactions numérotées « n° 371 » et « n° 392 »), au dos desquels il a
parfois écrit.\footnote{Pages blanches : 10, 18, 21. Tapuscrits sans trace de
sa main, non lus ici : 7, 9, 12, 20, les pages impaires de 23 à 39, et 41, qui
forme à elle seule le dernier lot du dossier. La page 17 ne porte que les deux
mots « Façon d'exprimer ».} Il est ordonné à l'envers de sa logique, et il faut
le savoir avant d'y entrer.

\begin{itemize}
\item \textbf{Le théorème et sa démonstration} — pages 1 à 5 : une première
liste de huit conditions sur un foncteur $F$, puis la démonstration, par
récurrence sur la dimension de la base, qu'elles suffisent pour que $F$ soit
représentable par un schéma localement quasi fini et séparé.
\item \textbf{Des reformulations} — pages 6 et 8 (la page 6 reformule la
page 8, et lui est donc postérieure), puis leur preuve à la page 19 : les
conditions de modularité réécrites avec des épaississements nilpotents et des
immersions ouvertes.
\item \textbf{Le cas non ramifié} — page 11 : une seconde liste, qui est le
brouillon du Corollaire 2 de la page 15.
\item \textbf{L'énoncé au net} — pages 13 à 16 : le Théorème, deux
corollaires (le cas étale, le cas non ramifié), deux remarques, des exemples
et des « applications en vue ».
\item \textbf{Les outils} — pages 22 à 40 : la Proposition 1 et ses deux
corollaires (un foncteur est représentable dès qu'il l'est sur les anneaux
locaux, puis sur les anneaux artiniens), la Proposition 2 (recoller le point
fermé d'une base locale complète au reste), la Proposition 3 (passer des
anneaux locaux de la base à la base entière), et la notion de paire
\emph{modulaire}.
\end{itemize}

Que les pages 22 à 40 soient les « notes précédentes » que citent les pages 3
à 5, les numéros le disent : la démonstration de la page 3 renvoie au
« cor 2 à prop. 1 », qui est le Corollaire 2 de la page 28 ; à la « prop. 2
des notes », qui est la Proposition 2 des pages 30 à 32 ; à la « prop. 3 des
notes préc. » et à sa « condition de constructibilité », qui sont la Prop. 3
de la page 32.\footnote{L'identification est la nôtre : les pages ne disent
pas où se trouvent les notes qu'elles citent. Mais les trois renvois tombent
chacun sur l'énoncé qui fait exactement ce que la démonstration lui demande.}
Et les deux programmes notés en marge de la page 32 — « transformer 4°) en une
condition plus facile à vérifier », avec une condition de dimension 1, et
« transformer en une condition de modularité générique » — sont réalisés par
le « critère de modularité a) locale » et « b) générique » de la page 1. Les
conditions (2) et (4) du Théorème de la page 13 sont, de même, écrites pour
nourrir le Corollaire 1 de la page 24 : elles ne sont demandées que pour
$B \to \widehat{B}$ et $\widehat{B} = \varprojlim B_n$, qui sont les deux
flèches de ce corollaire.

La datation est celle de l'inventaire, « vers 1971 » ; le dossier est rangé
dans le groupe « Descentes », que l'inventaire date « avant 1970 ». Les deux
ne s'accordent pas, et rien dans les pages ne tranche.

\subsection*{Les conventions}

\emph{Schémas et foncteurs.} Le manuscrit dit « préschéma » ; on écrit
« schéma ». $S$ est un schéma localement noethérien, $\mathbf{Sch}_{/S}$ la
catégorie des $S$-schémas, et $F : \mathbf{Sch}_{/S}^{\mathrm{op}} \to
\mathbf{Ens}$ un foncteur. Pour un $S$-schéma $S'$, $F_{S'}$ est la
restriction de $F$ aux $S'$-schémas. Pour $a \in S$, on pose $S_a =
\operatorname{Spec} \mathcal{O}_{S,a}$ et $F_a = F_{S_a}$. On identifie un
anneau $A$ au-dessus de $S$ à $\operatorname{Spec} A$, et l'on écrit $F(A)$.

\emph{Les numéros des conditions.} Les conditions du Théorème sont numérotées
(1) à (8) comme à la page 13. La liste de la page 1 est numérotée autrement ;
la concordance est la suivante :
\begin{itemize}
\item page 1, 1° a) et b) : conditions (1) et (2) ;
\item page 1, 2° et 3° : conditions (3) et (4) ;
\item page 1, 4° : sans équivalent à la page 13 ; on l'appelle (E) ;
\item page 1, 5° à 8° : conditions (5) à (8).
\end{itemize}
Aux pages 22 à 40, les conditions portent d'autres noms, qu'on garde :
(L1) à (L4), (R$_{\mathrm{loc}}$), (R$_0$).\footnote{Il trace sous chaque
$L$ une flèche, tournée vers la gauche, sauf pour $L_4$ aux pages 24 et 28 où
elle est tournée vers la droite, et marque $L_3$, $L_4$ de deux petits
cercles. La lecture ne leur prête aucun sens et les omet. (L1) est la
condition (1), (L2) la condition (3), (L3) et (L4) des formes faibles
(injectivité seulement) de (2) et (4).}

\emph{Modulaire.} C'est la notion centrale du dossier ; elle n'est définie
qu'à la page 34, mais elle est employée dès la page 1, et on la fixe ici. Soit
$T$ un $S$-schéma local, de point fermé $t$ au-dessus de $s \in S$, et $\eta
\in F(T)$. On dit que $(T, \eta)$ est \emph{modulaire} si, en gros, $\eta$ est
l'élément universel de $F$ au voisinage du point qu'il définit : tout
élément de $F$ sur un schéma local qui « passe par le même point » provient de
$\eta$ par un unique morphisme local (définition exacte à la page 34, plus
bas). Pour $\xi \in F(Z)$ et $z \in Z$, on dit que $\xi$ est modulaire en
$z$ si $(\operatorname{Spec} \mathcal{O}_{Z,z}, \xi_z)$ l'est. Lorsque $F_a$
est représentable par un schéma $X_a$ localement de type fini, $\xi$ définit
$u : Z \to X_a$ près de $z$, et $\xi$ est modulaire en $z$ si et seulement si
$\mathcal{O}_{X_a, u(z)} \to \mathcal{O}_{Z,z}$ est un isomorphisme : $u$ est
un isomorphisme local en $z$. C'est plus fort qu'« étale en $z$ », qu'il
écrit souvent entre crochets comme forme affaiblie (« modulaire [ou tout au
moins, étale] »).\footnote{Un morphisme étale en $z$ à extension résiduelle
triviale n'est pas pour autant un isomorphisme local : un revêtement étale de
degré 2 de la droite affine a des points rationnels au-dessus de points
rationnels. Les deux notions coïncident pour un monomorphisme, et aussi, à la
page 6, pour une section d'un foncteur non ramifié sur un sous-schéma fermé :
c'est là qu'il écrit « modulaire [i.e. étale] », à bon droit.}
\emph{Pro-modulaire}, qu'il emploie aux pages 4, 8 et 13 sans le définir, est
lu ici comme la même condition restreinte aux $T'$ artiniens : le complété de
$\mathcal{O}_{Z,z}$ pro-représente $F$ au point considéré. Enfin « étale »,
pour $\xi : Z \to F$ lorsque $F$ n'est pas représentable, s'entend au sens de
le relèvement infinitésimal (formellement étale et localement de
présentation finie).

\subsection*{Ce que le dossier annonce et n'établit pas}

\begin{itemize}
\item la suffisance des conditions du Théorème repose sur la démonstration
des pages 3 à 5, dont la dernière étape s'arrête sur « on est ramené à
vérifier la condition de constructibilité » : la déduction de cette
condition à partir de (8) n'est pas écrite ;
\item à la même démonstration, l'effectivité de la donnée de descente de
l'étape (3) est affirmée sans argument ;
\item le Lemme de la page 38 est énoncé sans démonstration, et la
« proposition formelle » qui recolle les $X^{(x)}$ à la page 40 est nommée,
pas écrite ;
\item la page 14 (« Immersion ouverte »), la Remarque 4 de la page 16, la
condition (D) de la page 6 et la page 17 s'interrompent ;
\item l'« exemple type » réclamé en marge de la page 13, qui montrerait que
(8) ne découle pas des autres conditions, n'est pas donné ;
\item les deux « applications en vue » de la page 16 sont annoncées, et la
première est fausse telle qu'elle est écrite.
\end{itemize}

\pagerange{1}{2}
\section*{Une première liste de conditions (pages 1 et 2)}

La page 1, marquée d'un « A » cerclé, pose huit conditions sur $F$. On les
donne dans la forme où la suite les emploie.
\begin{enumerate}
\item[1°] $F$ est un faisceau pour la topologie fpqc, c'est-à-dire : a) $F$
est de nature locale pour la topologie de Zariski, et b) $F$ satisfait la
descente pour les morphismes fidèlement plats quasi compacts affines.
\item[2°] $F$ commute aux limites inductives filtrantes d'anneaux (propriété
relative à $S$) : $F$ est \emph{localement de présentation finie}.
\item[3°] $F$ commute aux limites projectives adiques d'anneaux artiniens :
pour $B$ local noethérien complet, $F(B) \to \varprojlim F(B/\mathfrak{m}^n)$
est bijective.
\item[4°] Exactitude (E) : pour tout anneau local artinien $A$ au-dessus de
$S$ et toute $A$-algèbre finie locale $A' \supset A$ de même corps résiduel,
avec $A'/A$ de longueur 1 comme $A$-module, la suite
$F(A) \to F(A') \rightrightarrows F(A' \otimes_A A')$ est exacte.
\item[5°] Pour tout corps $k$ au-dessus de $S$ et tout $\xi \in F(k)$, $F_k$
est pro-représentable en $\xi$, et l'anneau qui le pro-représente est
\emph{fini} sur $k$.
\item[6°] $F$ est séparé sur $S$, au sens du critère valuatif.
\item[7°] Critère de modularité local : pour tout $S$-schéma local
noethérien complet $S'$, tout $S'$-schéma local fini $Z$ de point fermé $t$,
et tout $\xi \in F(Z)$ modulaire pour $F_{S'}$ en $t$, $\xi$ est modulaire
(ou au moins étale) en tout point de $Z$.
\item[8°] Critère de modularité générique : pour tout $S$-schéma noethérien
irréductible $S'$, tout $S'$-schéma fini $Z$ et tout $\xi \in F(Z)$ modulaire
pour $F_{S'}$ en un point $z$ au-dessus du point générique de $S'$, il existe
un voisinage ouvert $U$ de $z$ dans $Z$ en tout point duquel $\xi$ est
modulaire (ou au moins étale).
\end{enumerate}

Les quatre premières sont réunies sur la page par une accolade et une mention
au crayon, « conditions d'exactitude » ; deux notes marginales qualifient 5°
de condition de « finitude ou discrétion » et 6° de « séparation ». Ce sont
toutes des conditions \emph{nécessaires} : un foncteur représentable par un
schéma localement quasi fini et séparé les satisfait.\footnote{Pour (E), qui
est la seule dont ce ne soit pas immédiat, l'argument, que la page ne donne
pas, est le suivant. Un morphisme de $\operatorname{Spec} A'$ dans $X$ tombe
dans un ouvert affine, et l'on est ramené à voir que $A$ est l'égalisateur de
$A' \rightrightarrows A' \otimes_A A'$. Si $a' \in A' \smallsetminus A$
vérifiait $a' \otimes 1 = 1 \otimes a'$, l'image de ces deux éléments par
$\mathrm{id} \otimes \pi : A' \otimes_A A' \to A' \otimes_A (A'/A)$ serait
$0$ d'un côté, et de l'autre la classe de $1$ dans $A'/\mathfrak{m}_A A'$,
puisque $A'/A \simeq A/\mathfrak{m}_A$ est engendré par l'image de $a'$ ; or
cette classe n'est pas nulle, par Nakayama. La page écrit la condition avec
« ext.\ résiduelles triviales » et « long$_A A'/A = 1$ » entre crochets ; la
lecture en fait les hypothèses.} Pour 5°, le mot « corps » est une correction
au-dessus d'un premier « anneau local » biffé. Pour 7°, une note marginale
ajoute qu'on peut se borner à $S'$ irréductible de dimension 1, et un petit
diagramme place $\xi$ au-dessus de $F_{S'} \to F$ :

\begin{tikzcd}
F \arrow[d] & F_{S'} \arrow[l] \arrow[d] & \\
S & S' \arrow[l, no head] & Z \arrow[l] \arrow[ul, "\xi"']
\end{tikzcd}

Une note au crayon, en travers de la ligne de 5°, sur $F(k) \to F(X)$, est
presque entièrement illisible et n'est pas lue ici.

La page 2, au crayon, indique ce qui change pour les foncteurs non ramifiés :
remplacer 5° par la non-ramification de $F$, et simplifier 4° et 5° en
conséquence. C'est ce que fera le Corollaire 2 de la page 15.

\pagerange{3}{5}
\section*{La démonstration par récurrence sur la dimension (pages 3 à 5)}

Les conditions sont stables par changement de base, et toutes, sauf (3), (5)
et (6), sont de nature locale sur $F$.\footnote{La page dit « sauf 2°, 5°,
6° » dans la numérotation de la page 1 ; c'est (3), (5), (6) dans celle du
Théorème. Dans cette section, tous les numéros de la page sont traduits
selon la concordance donnée plus haut.} La démonstration traite d'abord le
cas $\dim S < +\infty$, par récurrence, en quatre étapes.

\emph{(1) $\dim S = 0$.} On se ramène à $S$ local artinien. Grâce à (E) et
(5), on construit un schéma $X$ localement quasi fini sur $S$ — une somme de
spectres d'anneaux artiniens finis sur $S$ — et, grâce à (1) à (4), il
représente $F$. C'est le Corollaire 2 de la page 28, qui est cité (« cf.\
notes précédentes, cor 2 à prop.\ 1 »). Seules les conditions (1) à (5) et
(E) ont servi.

\emph{(2) $\dim S = n > 0$, $S$ local complet.} Soit $s$ le point fermé et
$U = S \smallsetminus \{s\}$. Par l'hypothèse de récurrence, $F_U$ est
représentable par un $X_U$.\footnote{Entre crochets : « où 1° à 5° suffisent
si $S$ de dim 1 ». Une note marginale ajoute que si $\dim S = 1$, l'hypothèse
(8) est inutile.} Grâce à (3) et (5), et parce que $S$ est complet, on
construit une famille de $S$-schémas \emph{finis} $X_i =
\operatorname{Spec} B_i$, un par point de $F_s$ ; grâce à (4), des morphismes
$u_i : X_i \to F$. La Proposition 2 des pages 30 à 32 ramène alors à
vérifier que les restrictions $u_i|_U : X_{i,U} \to X_U$ sont des immersions
ouvertes, d'images deux à deux disjointes.

C'est là qu'interviennent (6) et (7). Par (7), il suffit que $u_i$ soit
modulaire au point fermé $x_i$ de $X_i$. Sa restriction à la fibre fermée
l'est, de sorte que le corps $k(x_i)$ est un corps de définition du point
qu'il définit, et il reste la condition de modularité proprement dite. Il la
contourne en démontrant d'abord le théorème sous une forme renforcée de (7),
où l'on suppose seulement $\xi$ \emph{pro-modulaire} en $z$ : c'est le cas de
$u_i$ par construction, puisque $B_i$ pro-représente $F$.\footnote{Une note
marginale dit au contraire qu'on n'utilise « pour l'instant » qu'« une
version affaiblie » de 7° ; la lecture de ces mots est douteuse. Comme
condition sur $F$, exiger la conclusion de (7) sous l'hypothèse plus faible
de pro-modularité est une condition plus forte. Elle reste nécessaire : si $X$
est localement quasi fini sur $S'$ local complet, donc hensélien,
l'anneau local $\mathcal{O}_{X,x}$ en un point au-dessus du point fermé est
fini sur $\mathcal{O}_{S'}$ (théorème principal de Zariski), donc complet, et
pro-modulaire y équivaut à modulaire. Le théorème n'est donc pas affecté.} Le
reste de l'étape — que les $u_i|_U$ soient des monomorphismes d'images
disjointes, d'où des immersions ouvertes puisqu'ils sont étales — se fait par
l'argument de la page 19 (plus bas) : le produit fibré $X_i \times_F X_j$ est
un sous-schéma fermé de $X_i \times_S X_j$, fini sur le complet $S$, qui
coïncide formellement avec la diagonale si $i = j$ et est formellement vide si
$i \neq j$ ; il est donc égal à la diagonale, ou vide.\footnote{La page
s'arrête sur « Ce dernier point provient de l'hypothèse 6° » et « utilisant
la condition 7° » ; l'enchaînement par $X_i \times_F X_j$ est emprunté à la
page 19, qui fait ce raisonnement pour l'axiome (A). Le « 6° » entouré en tête
de la page 4 achève la phrase de la page 3 et ne désigne pas une étape.} Enfin
$X$ est localement quasi fini sur $S$ par construction, et séparé par (6).

\emph{(3) $\dim S = n$, $S$ local.} Par l'étape (2), $F_{\widehat{S}}$ est
représentable par un $X_{\widehat{S}}$ localement quasi fini et séparé sur
le complété $\widehat{S}$. La donnée de descente canonique relative à
$\widehat{S} \to S$ est effective, d'où un $X$ sur $S$, localement quasi fini
et séparé ; par (2), il représente $F$.\footnote{La page affirme l'effectivité
sans argument. Sur une partie quasi compacte, un schéma séparé et quasi fini
est quasi affine (théorème principal de Zariski), et la descente fidèlement
plate des schémas quasi affines est effective (SGA 1, VIII). Pour un
$X_{\widehat{S}}$ qui n'est pas quasi compact, il faudrait le recouvrir par
des ouverts quasi compacts stables par la donnée de descente ; la page ne le
fait pas, et la présente lecture non plus.}

\emph{(4) $S$ quelconque.} On utilise le

\textbf{Lemme.} Si $F_{S_a}$ est représentable pour tout $a \in S$, alors
$F$ est représentable.

La page 4 renvoie pour cela à la Proposition 3 de la page 32, et la page 5,
une seule ligne au crayon, conclut : compte tenu de (1) et (3), on est ramené
à vérifier la condition de constructibilité de cette proposition. Le dossier
s'arrête là. La condition (8), le critère de modularité générique, est faite
pour fournir cette constructibilité — la note marginale de la page 32 le dit
en toutes lettres —, mais la déduction n'est écrite nulle part.\footnote{Le
Lemme tel qu'il est énoncé n'est vrai que sous les autres conditions du
théorème (il s'agit du « $F$ comme dessus ») ; ce sont celles de la
Proposition 3 qu'il utilise.}

\pagerange{6}{8}
\section*{Épaississements et immersions ouvertes (pages 6 et 8)}

Les deux pages reformulent les critères de modularité. La page 8 est la
première : la page 6 renvoie à ses axiomes (A) et (B).

\subsection*{Les axiomes (A), (B), (C) (page 8)}

Il garde les conditions (1) à (6) et remplace (7) et (8) par :

\textbf{(A)} Si $S'$ est local noethérien complet (de dimension 1 si l'on
veut), $X$ fini sur $S'$, et $\xi : X \to F_{S'}$ pro-modulaire en un point
$z$ au-dessus du point fermé $s'$ de $S'$, alors $\xi$ est une immersion
ouverte — pourvu que $F$ soit représentable au-dessus de $S' \smallsetminus
\{s'\}$.

\textbf{(B)} Si $S'$ est noethérien sur $S$, $\xi : X \to F_{S'}$ et $x \in
X$ tels que $\xi$ soit modulaire en $x$, il existe un voisinage ouvert $U$ de
$x$ tel que $\xi|_U : U \to F_{S'}$ soit une immersion ouverte.\footnote{Entre
crochets, il précise : « on veut dire simplement » que la condition porte sur
le changement de base à $\operatorname{Spec} \mathcal{O}_{S',s'}$. Le mot qui
suit « $X$ » dans l'hypothèse est illisible ; on lit $X$ de type fini sur
$S'$, ce qu'exige la conclusion.}

Pour $F \to S$ un monomorphisme, (A) et (B) contiennent la condition (C) :
pour tout $S'$ noethérien sur $S$ et tout $\xi : X \to F_{S'}$, l'ensemble
des points de $X$ où $\xi$ est étale est ouvert.\footnote{Un long trait
oblique traverse ce passage, du « (i) » biffé à « l'ens des » ; on ne sait
s'il le biffe. La lecture le garde, parce que la suite de la page le
reprend.} Un lemme permet de supposer, dans (A) et (B), que $X \to S'$ est une
immersion fermée, définie si l'on veut par un idéal de carré nul. Si $F \to S$
est non ramifié, (A) et (B) se réduisent donc à :

\textbf{(C)} Pour tout $S'$ noethérien sur $S$, toute immersion fermée $S'_0
\to S'$ et tout $\xi : S'_0 \to F_{S'}$, l'ensemble $U$ des points où $\xi$
est étale (et injectif) est ouvert ; c'est-à-dire :
\begin{itemize}
\item[(C$_1$)] $U$ est stable par générisation (condition équivalente à (A),
dit-il) ;
\item[(C$_2$)] pour tout $x \in U$, $U \cap \overline{\{x\}}$ est un
voisinage de $x$ dans $\overline{\{x\}}$.
\end{itemize}
L'équivalence est exacte dans un espace noethérien.\footnote{Si $U$ vérifie
(C$_1$) et (C$_2$), son complémentaire $V$ est stable par spécialisation ;
le point générique d'une composante irréductible de $\overline{V}$ ne peut
être dans $U$, sans quoi (C$_2$) enlèverait à cette composante un ouvert non
vide ; donc $\overline{V}$ est réunion des adhérences de points de $V$, et
$V$ est fermé. La page dit « inj.\ étale » ; on lit « étale et injectif ».}
La condition (C) est automatique si $F \to S$ est étale.

\subsection*{La reformulation de la page 6}

La page 6 réécrit (C$_1$) comme un critère d'extension, sous la forme d'une
nouvelle condition 7° :

\emph{Pour tout $S$-schéma $S'$ local noethérien complet irréductible de
dimension 1, d'ouvert $U = S' \smallsetminus \{\text{point fermé}\}$
schématiquement dense dans $S'$, pour tout sous-schéma $S'_0$ défini par un
idéal nilpotent (de carré nul), et tout $S'$-morphisme $\xi_0 : S'_0 \to
F_{S'}$ : si $\xi_0|_{U_0}$ se prolonge en $\xi_U : U \to F_{S'}$, alors
$\xi_0$ se prolonge en $\xi : S' \to F_{S'}$.}

Elle est vide si $F \to S$ est étale : le relèvement le long d'un
épaississement nilpotent existe alors toujours. Et (C$_2$) se met sous la
forme d'une nouvelle 8° : si $S'$ est noethérien irréductible de point
générique $x$, $S'_0$ un sous-schéma fermé contenant $x$ et $\xi : S'_0 \to
F_{S'}$ modulaire — c'est-à-dire ici étale — en $x$, alors $\xi$ est
modulaire sur un voisinage de $x$.\footnote{« d'adhérence schématique $S'$ »
est une note entourée dans la marge, reliée à « $U \subset S'$ » ; la
lecture en fait l'hypothèse de densité schématique. Elle est nécessaire au
sens de la condition : sans elle, $U_0 = U$ dès que l'idéal de $S'_0$ est à
support dans le point fermé, et la condition exigerait qu'on relève tout, ce
que ne fait pas même une immersion fermée (voir le Corollaire 2, page 15).}

Suit, « revenant au cas général (pas nécessairement non ramifié) », une
condition (D) qui devait remplacer (A) et (B) : pour $S'$ noethérien sur $S$,
$X$ fini sur $S'$ et $\xi : X \to F_{S'}$, l'ensemble des points où $\xi$ est
étale est ouvert, et de même pour $X \times_{F_{S'}} X$ sur $S'$, via la
diagonale $F \to F \times_S F$. Le bas de la page est un palimpseste où
seuls quelques mots se lisent, et la condition n'est pas achevée.

\pagerange{11}{11}
\section*{Le cas non ramifié : une seconde liste (page 11)}

La page 11 est le brouillon du Corollaire 2 de la page 15. Elle reprend (1) à
(6) (la nature locale, la descente fidèlement plate quasi compacte, la
commutation aux limites inductives, aux limites projectives d'anneaux
artiniens, la pro-représentabilité universelle, la séparation), remplace (7)
par « $F$ non ramifié », et (8) par une condition d'extension sur les anneaux
locaux de dimension 1, qu'on retrouvera sous sa forme juste à la page 15. Une
condition (6°) sur les points à valeurs dans une extension $K$, qui devaient
provenir d'une sous-extension finie, est biffée.

Elle ajoute une condition nouvelle, numérotée (9°) :

\textbf{(RG)} Pour tout $T$ noethérien irréductible sur $S$, il existe un
ouvert non vide $U \subset T$ tel que $F_U$ soit représentable.

C'est la condition de représentabilité générique que la page 16 nommera
(RG). Une note marginale indique qu'elle peut se formuler plus faiblement,
par une impossibilité de prolongement. Deux colonnes de notes verticales,
dans la marge gauche, ne se lisent pas.

\pagerange{13}{13}
\section*{Le théorème (pages 13 à 16)}

\subsection*{L'énoncé}

\textbf{Théorème.} Soient $S$ un schéma localement noethérien et $F :
\mathbf{Sch}_{/S}^{\mathrm{op}} \to \mathbf{Ens}$. Pour que $F$ soit
représentable par un $S$-schéma $X$ localement quasi fini et séparé sur $S$,
il faut et il suffit que $F$ satisfasse les conditions suivantes :
\begin{enumerate}
\item[(1)] $F$ est de nature locale (un faisceau Zariski) ;
\item[(2)] $F$ est compatible à la descente fidèlement plate quasi compacte ;
il suffit de le demander pour $B \to \widehat{B}$, où $B$ est un anneau local
noethérien au-dessus de $S$ ;
\item[(3)] $F$ commute aux limites inductives filtrantes d'anneaux ;
\item[(4)] $F$ commute aux limites projectives d'anneaux artiniens de la forme
$\widehat{B} = \varprojlim B_n$, $B$ comme dans (2) ;
\item[(5)] au-dessus de tout anneau artinien $A$ sur $S$, $F$ est
pro-représentable, et les anneaux qui le pro-représentent sont finis sur $A$ ;
\item[(6)] $F$ est séparé (critère valuatif) ;
\item[(7)] pour tout anneau local noethérien complet $A$ sur $S$, tout
homomorphisme local fini $A \to B$ et tout $\xi \in F(B)$ modulaire au point
fermé, $\xi_{\mathfrak{p}} \in F(B_{\mathfrak{p}})$ est modulaire pour tout
$\mathfrak{p} \in \operatorname{Spec} B$ ;
\item[(8)] pour tout $S$-schéma noethérien irréductible $S'$, tout
$S'$-schéma fini $Z$, tout $\xi \in F(Z)$ et tout point $z \in Z$ au-dessus du
point générique de $S'$ en lequel $\xi$ est modulaire, il existe un voisinage
ouvert $U$ de $z$ dans $Z$ tel que $\xi$ soit modulaire (ou simplement étale)
en tout point de $U$.
\end{enumerate}

Sur la page, les trois premières conditions viennent dans l'ordre 1°, 3°,
2° ; la dernière commence à la page 13 et s'achève en tête de la page
15.\footnote{La page 15 écrit « en $Z$ » là où l'on attend le point $z$, et
« une partie ouverte non vide $U$ de $Z$ ». Comme $Z$ est fini sur $S'$, un
point au-dessus du point générique de $S'$ est générique dans $Z$, et, $Z$
étant irréductible comme le précise la page 1, un ouvert non vide le
contient : les deux formes disent la même chose. On écrit la plus claire.}
La nécessité se vérifie condition par condition.\footnote{Pour (7) : si
$\xi : \operatorname{Spec} B \to X$ induit $\mathcal{O}_{X,x} \simeq B$ au
point fermé, c'est le morphisme canonique $\operatorname{Spec}
\mathcal{O}_{X,x} \to X$, qui est un isomorphisme local en chaque point. Pour
(8) : un isomorphisme d'anneaux locaux entre schémas localement de
présentation finie sur $S'$ se prolonge en un isomorphisme de voisinages
(EGA IV, 8). La page n'écrit pas ces vérifications.} La suffisance est ce que
démontrent les pages 3 à 5, avec les lacunes dites plus haut.

Il ajoute trois remarques. D'abord, dans (7), l'hypothèse porte sur $\xi$
modulaire, « et l'on verra que cela implique pro-modulaire » ; en marge : « en
fait, cela suffit, et cela signifie que $\xi : \operatorname{Spec} B \to F$
est étale ». Ensuite, qu'on pourrait montrer que (1) à (6) et (7) impliquent
l'analogue de (7) sans l'hypothèse de dimension 1 — ce qui confirme que (7)
était pensée d'abord sur une base de dimension 1, comme le disait la marge de
la page 1. Enfin, une question : « Donner l'exemple type [$S$ de dim 1]
montrant que (8°) n'est pas conséquence des autres. Est-ce que 6° et 7°
[suffisent] ? \emph{Non} ». L'exemple n'est pas donné.\footnote{Deux notes
obliques, dans la marge gauche, dont plusieurs mots sont illisibles ; les
crochets de cette phrase sont nôtres.}

\pagerange{14}{14}
\subsection*{Immersion ouverte (page 14)}

Sur un feuillet isolé, le début d'un corollaire pour les immersions
ouvertes : sous (1) à (4), et si de plus $F \to S$ est un monomorphisme étale
pour les arguments artiniens, $F$ devrait être représentable par un ouvert de
$S$. La conclusion n'est pas écrite.

\pagerange{15}{16}
\subsection*{Les deux corollaires (page 15)}

\textbf{Corollaire 1.} Pour que $F$ soit représentable par un $S$-schéma
\emph{étale} et séparé, il faut et il suffit qu'il satisfasse (1), (2), (3),
(4), (6), et :

(Ét) $F$ est étale (ou seulement étale pour les arguments artiniens).

La liste omet (5), (7) et (8) : qu'elles découlent alors des autres, la page
ne le démontre pas, et une note marginale dit seulement que ce cas est « plus
direct et plus facile ».

\textbf{Corollaire 2.} Pour que $F$ soit représentable par un $S$-schéma
\emph{non ramifié} et séparé, il faut et il suffit que $F$ satisfasse (1) à
(6), et la forme simplifiée suivante de (7) :
\begin{itemize}
\item[a)] $F$ est non ramifié (les anneaux qui le pro-représentent dans (5)
sont non ramifiés sur $A$) ;
\item[b)] pour tout anneau local noethérien $A$ de dimension 1 sur $S$, de
spectre irréductible, d'idéal premier minimal $\mathfrak{p}$, tel que $A \to
A_{\mathfrak{p}}$ soit injectif, pour tout idéal nilpotent $I$ de $A$ et tout
$\eta \in F(A/I)$ : si l'image $\eta_{\mathfrak{p}}$ de $\eta$ dans
$F(A_{\mathfrak{p}}/I_{\mathfrak{p}})$ provient de $F(A_{\mathfrak{p}})$,
alors $\eta$ provient de $F(A)$.
\end{itemize}

La condition b) est corrigée ici.\footnote{La page porte, entre crochets,
« [idéal] $I$ [tel que $I_{\mathfrak{p}} = 0$] », sans hypothèse sur les
idéaux associés de $A$. Sous cette forme, $\eta_{\mathfrak{p}}$ provient
toujours de $F(A_{\mathfrak{p}})$, et la condition exige que $F(A) \to
F(A/I)$ soit surjective pour tout $I$ nilpotent nul en $\mathfrak{p}$ : c'est
faux déjà pour une immersion fermée. Soit $A = k[[t, \varepsilon]]/
(\varepsilon^2, t\varepsilon)$, de dimension 1, de spectre irréductible, de
premier minimal $\mathfrak{p} = (\varepsilon)$, et $I = (\varepsilon)$, qui
est nul dans $A_{\mathfrak{p}}$ puisque $t$ y est inversible ; soit $S =
\operatorname{Spec} A$ et $F$ le foncteur représenté par le sous-schéma fermé
$\operatorname{Spec} A/I$ de $S$, qui est non ramifié et séparé. Alors
$F(A/I)$ n'est pas vide, et $F(A)$ l'est. La lecture suit la page 11 (« tout
idéal $I$ de $A$ nilpotent ») et la note marginale de la page 6
(« d'adhérence schématique ») : $A \to A_{\mathfrak{p}}$ injectif, $I$
nilpotent quelconque. Sous cette forme la condition est nécessaire : un
schéma non ramifié est localement fermé dans un schéma étale, le relèvement
dans le schéma étale existe et est unique le long de $A \to A/I$, et il se
factorise par le fermé parce qu'il le fait génériquement et que $A \to
A_{\mathfrak{p}}$ est injectif. Que $\mathfrak{p}$ soit le premier minimal
est aussi une lecture : la page ne dit pas quel idéal premier.} Une note
marginale veut qu'on explique ici la condition (8) par une impossibilité de
prolongement ; c'est ce que fait b).

\subsection*{Deux remarques (pages 15 et 16)}

\textbf{Remarque 3.} La condition (8) est conséquence des autres, pour autant
que $F$ satisfasse :

\textbf{(RG)} pour tout sous-schéma fermé irréductible réduit $S'$ de $S$, il
existe un ouvert non vide $U \subset S'$ tel que $F_U$ soit représentable par
un schéma localement de type fini sur $U$.\footnote{La phrase de la Remarque 3
est en partie illisible ; que (RG) en soit l'hypothèse est une reconstruction,
appuyée sur la place de (RG) en tête de la page suivante et sur l'argument de
nécessité de (8), qui ne demande que la représentabilité au voisinage du point
générique.}

\textbf{Remarque 4.} Il faudrait limiter les anneaux $A$, $B$ et les schémas
$S'$ qui doivent être envisagés dans les conditions (2), (4), (5), (7), (8).
La remarque s'arrête sur un tiret.

\pagerange{16}{17}
\subsection*{Exemples et applications en vue (page 16)}

Il dresse la liste des foncteurs non ramifiés auxquels le Corollaire 2 doit
s'appliquer :
\begin{itemize}
\item les quotients $G/H$ d'un schéma en groupes par un sous-groupe fermé
$H$, sous une hypothèse qui les rende non ramifiés ;\footnote{Il avait écrit
« ouvert et fermé », puis biffé « ouvert ». Le premier mot de la ligne,
qualifié de « triviaux », est illisible. $G/H$ n'est pas non ramifié pour tout
$H$ fermé ; il l'est, par exemple, si $H$ est ouvert et fermé : la relation
d'équivalence est alors ouverte dans $G \times_S G$, et la diagonale de $G/H$
est une immersion ouverte.}
\item les classes de correspondances divisorielles, sous une hypothèse de
platitude ;
\item les foncteurs $\operatorname{Hom}$ à valeurs dans un groupe
commutatif, et en particulier les $\operatorname{Hom}$ de schémas abéliens,
non ramifiés par rigidité ;
\item les immersions, qui s'introduisent dans les problèmes de modules par
exemple via les isogénies, et en particulier la diagonale $F \to F \times_S
F$ ;
\item les sous-schémas en groupes d'un schéma en groupes ;
\item la décomposition des fibres d'un morphisme en composantes connexes (ou
irréductibles), qui donne même des foncteurs étales ;
\item la factorisation d'un morphisme fini.
\end{itemize}

Puis deux « applications en vue ». La première : un schéma abélien sur une
base localement noethérienne connexe serait projectif, en utilisant le schéma
des correspondances divisorielles sur $X \times_S X$. Sous cette forme,
l'énoncé est faux.\footnote{Raynaud (\emph{Faisceaux amples sur les schémas
en groupes}, 1970) a construit des schémas abéliens qui ne sont pas
projectifs sur leur base ; un schéma abélien est projectif localement pour la
topologie de Zariski lorsque la base est normale. Référence citée de mémoire.
La page ne cite pas Raynaud, et l'on ne dit rien ici de l'ordre des deux
textes. Le mot « connexe » est d'une lecture douteuse.} La seconde :
l'existence du schéma de Picard d'un schéma propre sur un corps, attribuée
entre crochets à « Chevalley-Seshadri-Grothendieck-Murre », et plus
généralement, sur une base $S$, son existence au-dessus d'un ouvert non vide.
La phrase s'arrête sur un tiret.\footnote{Le théorème d'existence sur un
corps est aujourd'hui attribué à Murre (1964). Les « correspondances
divisorielles » de ces deux lignes sont les points d'un foncteur non ramifié
de la liste ci-dessus, ce qui explique leur place ici.}

La page 17 ne porte que « Façon d'exprimer », et rien ne dit ce qui devait
suivre.

\pagerange{19}{19}
\section*{La preuve de (A) et (B) (page 19)}

La page 19 démontre les axiomes de la page 8. On étudie d'abord la diagonale
$F \to F \times_S F$ : elle est représentable, et, grâce à (6), c'est une
immersion fermée.\footnote{L'argument de représentabilité de la diagonale
est illisible sur la page ; la lecture ne le reconstruit pas.} Alors, dans
(A), $X \times_{F_{S'}} X$ est un sous-schéma \emph{fermé} de $X
\times_{S'} X$, fini sur $X$, qui coïncide formellement avec la diagonale
$X \to X \times_{S'} X$ — c'est ce que dit la pro-modularité — et lui est
donc égal, puisque tout est fini sur la base locale complète $S'$ et qu'un
schéma fini sur un anneau local complet est la somme de ses complétés en ses
points fermés. Donc $X \to F$ est un monomorphisme ; il est d'autre part
étale ; c'est une immersion ouverte.\footnote{La page écrit $X \times_S X$ ;
$X$ étant un $S'$-schéma, c'est le produit au-dessus de $S'$ qui est en jeu.
L'argument que « formellement égal » entraîne « égal » est nôtre ; la page
dit seulement « donc lui est identique ».}

Pour (B), on prouve de façon analogue qu'il existe un voisinage ouvert $U$ de
$x$ tel que $\xi|_U$ soit un monomorphisme ; la fin de la phrase, qui devait
conclure par « étale », est illisible.

\pagerange{22}{24}
\section*{Les notes précédentes : critères formels (pages 22 à 40)}

Ces pages sont logiquement les premières du dossier. Elles établissent, sans
hypothèse de finitude sur $X$, qu'un morphisme $\xi : X \to F$ est un
isomorphisme dès qu'il l'est sur des arguments de plus en plus petits : les
schémas locaux, puis les schémas artiniens.

\subsection*{Proposition 1 : des arguments locaux au cas général (pages 22 et 24)}

\textbf{Proposition 1.} Soient $S$ un schéma, $F :
\mathbf{Sch}_{/S}^{\mathrm{op}} \to \mathbf{Ens}$, $X$ un $S$-schéma
localement de présentation finie et $\xi \in F(X)$. On suppose :
\begin{itemize}
\item[(L1)] $F$ est de nature locale ;
\item[(L2)] $F$ commute aux limites inductives filtrantes d'anneaux ;
\item[(R$_{\mathrm{loc}}$)] $(X, \xi)$ représente $F$ sur les $S$-schémas
locaux essentiellement de type fini sur $S$ dont le point fermé a une
extension résiduelle finie sur $S$.
\end{itemize}
Alors $(X, \xi)$ représente $F$ : $\xi : h_X \to F$ est un isomorphisme.

\emph{Démonstration.} Par (L1), on peut supposer $S$ affine ; par (L2), et
parce que $h_X$ commute aussi aux limites filtrantes ($X$ étant localement de
présentation finie), il suffit de montrer que $\operatorname{Hom}_S(T, X) \to
F(T)$ est bijective pour $T$ de présentation finie sur $S$. Appelons
\emph{admissibles} les points $t \in T$ dont l'extension résiduelle est finie
sur $S$ ; tout fermé non vide de $T$ en contient un (un point fermé de sa
fibre), de sorte qu'un ouvert contenant tous les points admissibles est
$T$ tout entier.

\emph{Injectivité.} Si $u, v : T \to X$ ont même image dans $F(T)$, ils ont
par (R$_{\mathrm{loc}}$) même restriction à $\operatorname{Spec}
\mathcal{O}_{T,t}$ pour tout $t$ admissible, donc — $X$ étant localement de
présentation finie — coïncident sur un voisinage de chaque tel $t$, donc
partout.\footnote{La page dit seulement « d'où la conclusion ». Le passage
par les points admissibles et le voisinage est nôtre.}

\emph{Surjectivité.} Soit $\eta \in F(T)$. Pour $t$ admissible, il existe par
(R$_{\mathrm{loc}}$) un $u_t : \operatorname{Spec} \mathcal{O}_{T,t} \to X$,
unique par ce qui précède, qui induit $\eta_t$. Comme $X$ est localement de
présentation finie, $u_t$ s'étend en $v_t : U_t \to X$ sur un voisinage
ouvert $U_t$ de $t$, et, par (L2), quitte à rétrécir $U_t$, on a
$F(v_t)(\xi) = \eta|_{U_t}$. Par l'injectivité, $v_t$ et $v_{t'}$ coïncident
sur $U_t \cap U_{t'}$ ; les $U_t$ recouvrent $T$ ; les $v_t$ se recollent en
$v : T \to X$, et par (L1), $F(v)(\xi) = \eta$.\footnote{La page renvoie à
« a) » et « b) », qui sont les premières étiquettes, biffées, de (L1) et
(L2).}

\pagerange{24}{28}
\subsection*{Corollaire 1 : des arguments artiniens aux arguments locaux (pages 24 à 28)}

\textbf{Corollaire 1.} Sous les hypothèses préliminaires de la Proposition 1,
supposons de plus $S$ localement noethérien, et, pour tout anneau local $B$
essentiellement de type fini sur $S$ à extension résiduelle finie, de
complété $\widehat{B} = \varprojlim B_n$ :
\begin{itemize}
\item[(L3)] $F(B) \to F(\widehat{B})$ est injective ;
\item[(L4)] $F(\widehat{B}) \to \varprojlim F(B_n)$ est injective ;
\item[(R$_0$)] $(X, \xi)$ représente $F$ sur les arguments artiniens
essentiellement de type fini sur $S$, à extension résiduelle finie.
\end{itemize}
Si (L1), (L2), (L3), (L4) et (R$_0$) sont satisfaites, $(X, \xi)$ représente
$F$.\footnote{(L3) est très raturée sur la page : posée d'abord sur un
complété, biffée ligne à ligne ; la note marginale en donne la forme retenue,
« $F(B) \to F(\widehat{B}) \to \varprojlim$, injectif ».}

\emph{Démonstration.} Il suffit de vérifier (R$_{\mathrm{loc}}$). Soit $T =
\operatorname{Spec} B$ comme ci-dessus, $T_n = \operatorname{Spec} B_n$, $T' =
\operatorname{Spec} \widehat{B}$, $f : T' \to T$, et $\eta' = F(f)(\eta)$.

\emph{Injectivité.} Si $F(u)(\xi) = F(v)(\xi)$, les restrictions $u_n, v_n$ à
$T_n$ ont même image, donc sont égales par (R$_0$) ; donc $u = v$, car $B$
est noethérien et $B \subset \widehat{B}$.

\emph{Surjectivité.} Par (R$_0$), il existe des $u_n : T_n \to X$ induisant
les $\eta_n$, uniques, donc compatibles, d'où $u' : T' \to X$ qui les induit,
et, par (L4), $F(u')(\xi) = \eta'$.\footnote{La page invoque ici une condition
dont l'indice est « loc » ; c'est (L4) qui donne l'égalité, deux éléments de
$F(\widehat{B})$ ayant même image dans $\varprojlim F(B_n)$.} Soient $p_1, p_2
: T' \times_T T' \rightrightarrows T'$ les deux projections :
\[
T' \times_T T' \;\rightrightarrows\; T' \xrightarrow{\ f\ } T,
\qquad u' : T' \to X .
\]
On a $F(u' p_1)(\xi) = F(p_1)(\eta') = F(p_2)(\eta') = F(u' p_2)(\xi)$,
puisque $f p_1 = f p_2$. Il en conclut $u' p_1 = u' p_2$ « en vertu de
1°) », c'est-à-dire de l'injectivité.\footnote{Telle qu'elle est prouvée,
l'injectivité ne vaut que pour des schémas locaux essentiellement de type
fini, et $T' \times_T T' = \operatorname{Spec}(\widehat{B} \otimes_B
\widehat{B})$ n'en est pas un. Elle s'étend pourtant : par l'argument de la
Proposition 1, elle vaut pour tout $T$ de présentation finie sur $S$, puis,
par (L2) et parce que $X$ est localement de présentation finie, pour tout
$S$-schéma affine, deux morphismes vers $X$ se factorisant par un même
$T_i$ de présentation finie dans une présentation en limite. La lecture
supplée cette extension, que la page ne fait pas.} Comme $f$ est fidèlement
plat, $u'$ descend : il existe $u : T \to X$ avec $u' = u f$. Enfin
$\eta = F(u)(\xi)$, car par (L3) il suffit de le vérifier après $F(f)$, où
cela s'écrit $\eta' = F(u')(\xi)$.

\emph{N.B.} Il note qu'il serait plus naturel d'énoncer ces deux résultats
pour un morphisme de foncteurs $G \to F$ satisfaisant les conditions
envisagées pour $F$, et de conclure que c'est un isomorphisme.

\pagerange{28}{30}
\subsection*{Corollaire 2 : une base artinienne (pages 28 et 30)}

\textbf{Corollaire 2.} Soit $S$ un schéma artinien. Pour que $F$ soit
représentable par un $X$ localement quasi fini sur $S$, il faut et il suffit
que :
\begin{itemize}
\item[1°)] $F$ soit pro-représentable, et que les anneaux locaux qui
interviennent dans la pro-représentation soient artiniens, c'est-à-dire finis
sur $S$ ;
\item[2°)] $F$ satisfasse (L1), (L2), (L3), (L4).
\end{itemize}

En effet, par 1°) et (L1), les anneaux de la pro-représentation donnent un
$X = \coprod_i \operatorname{Spec} B_i$ et un $\xi \in F(X)$ qui vérifient
(R$_0$) ; le Corollaire 1 conclut.\footnote{« Pro-représentable » au sens de
Grothendieck (\emph{Technique de descente II}, Séminaire Bourbaki 195) : sur
les anneaux artiniens, $F$ est somme de foncteurs pro-représentables par des
anneaux locaux, $F(A) = \coprod_i \operatorname{Hom}_{\mathrm{loc}}(B_i, A)$
pour $A$ local. La réciproque — qu'un $X$ localement quasi fini sur $S$
artinien est une somme de spectres d'anneaux artiniens finis sur $S$ — n'est
pas écrite.} C'est l'étape (1) de la démonstration de la page 3.

\pagerange{30}{32}
\subsection*{Proposition 2 : recoller le point fermé (pages 30 et 32)}

\textbf{Proposition 2.} Soient $A$ un anneau local noethérien complet, $S =
\operatorname{Spec} A$, $S' = S \smallsetminus \{\text{point fermé}\}$ et $F
: \mathbf{Sch}_{/S}^{\mathrm{op}} \to \mathbf{Ens}$. On suppose :
\begin{itemize}
\item[1°)] $F_{S'}$ est représentable par un schéma $X'$ localement de type
fini sur $S'$ ;
\item[2°)] la restriction de $F$ aux $A$-algèbres artiniennes finies est
pro-représentable, par des anneaux locaux \emph{finis} sur $A$ ;
\item[3°)] $F$ satisfait (L1), (L2), (L3), et (L4) sous la forme forte :
$F(\widehat{B}) \to \varprojlim F(B_n)$ est bijective.
\end{itemize}

Grâce à 2°) et 3°), on construit $X_1 = \coprod_i \operatorname{Spec} B_i$,
où les $B_i$ sont les anneaux de la pro-représentation, et $\xi_1 \in
F(X_1)$.\footnote{Pour passer du système compatible que donne la
pro-représentation à un élément de $F(B_i)$, il faut que $F(\widehat{B}) \to
\varprojlim F(B_n)$ soit \emph{surjective} : c'est l'effectivité des objets
formels, que (L4), telle qu'elle est écrite à la page 24, ne donne pas. La
lecture l'ajoute à l'hypothèse 3°). La condition (4) du Théorème de la page
13 la demande, puisqu'elle dit « commute ». Sur la page 30, la flèche sous le
quatrième $L$ paraît d'ailleurs tournée dans l'autre sens qu'aux pages 24 et
28. La page écrit aussi $X_1^{*}$, $\xi_1^{*}$, puis $X_1$ sans l'étoile ; on
n'en garde qu'une notation.} Soient $X'_1 = X_1 \times_S S'$ et $\xi'_1$ la
restriction de $\xi_1$ ; par 1°), $\xi'_1 = F(u)(\xi')$ pour un unique
morphisme $u : X'_1 \to X'$. On suppose enfin :
\begin{itemize}
\item[4°)] $u$ est une immersion ouverte.
\end{itemize}
Alors $F$ est représentable. En effet, on recolle $X_1$ et $X'$ le long de
$u$ ; on obtient $X$, et, par (L1), un $\xi \in F(X)$ qui induit $\xi_1$ et
$\xi'$. Le Corollaire 1 s'applique : un argument artinien d'extension
résiduelle finie est au-dessus du point fermé, où $X$ se réduit à $X_1$, ou
au-dessus de $S'$, où il se réduit à $X'$ ; (R$_0$) est donc satisfaite, et
$(X, \xi)$ représente $F$.

Deux notes marginales accompagnent la proposition. L'une, d'une autre plume,
la reformule : sont donnés $X_U \simeq F_U$ sur l'ouvert $U$ et $Z \to F$ sur
$S$, avec $Z_U \to X_U$ une immersion ouverte et $Z_0 \to F_0$ un
isomorphisme sur la fibre fermée. L'autre est un programme : « transformer
4°) en une condition plus facile à vérifier », avec une condition de
dimension 1. C'est le critère de modularité local (7).

\pagerange{32}{36}
\subsection*{Proposition 3 et les paires modulaires (pages 32 à 36)}

\textbf{Proposition 3.} Soit $F : \mathbf{Sch}_{/S}^{\mathrm{op}} \to
\mathbf{Ens}$. On suppose :
\begin{itemize}
\item[1°)] pour tout $a \in S$, $F_a$ est représentable par un schéma $X_a$
localement de type fini sur $S_a$ ;
\item[2°)] pour tout $S$-schéma $T$ localement de type fini et tout $\eta \in
F(T)$, l'ensemble $M(T, \eta)$ des $t \in T$ en lesquels $\eta$ est modulaire
est une partie \emph{constructible} de $T$ ;
\item[3°)] $F$ satisfait (L1) et (L2).
\end{itemize}
Alors $F$ est représentable.

En marge de 2°) : « Transformer en une condition de modularité
\emph{générique} ». C'est le critère (8).

\emph{Définition (page 34).} Soit $T$ un $S$-schéma local, de point fermé
$t$ au-dessus de $s \in S$, $k = k(s)$, et $\eta \in F(T)$, de valeur $\beta =
\eta(t) \in F(k(t))$ au point fermé. On dit que $(T, \eta)$ est
\emph{modulaire} si :
\begin{itemize}
\item[1°)] pour tout $S$-schéma local $T'$, de point fermé $t'$, l'application
$\operatorname{Hom}_{\mathrm{loc}, S}(T', T) \to F(T')$ définie par $\eta$
est une bijection sur l'ensemble des $\eta' \in F(T')$ dont la valeur
$\eta'(t') \in F(k(t'))$ est dans l'image de
\[
\operatorname{Hom}_k\bigl(k(t), k(t')\bigr) \longrightarrow F(k(t')),
\]
l'application étant définie par $\beta$ ;\footnote{La page écrit
$\operatorname{Hom}_k(k(t'), k(t))$ ; mais c'est un $k$-plongement de $k(t)$
dans $k(t')$ qui transporte $\beta \in F(k(t))$ dans $F(k(t'))$, et c'est la
variance qu'exige la phrase. La lecture la rétablit.}
\item[2°)] $L = k(t)$ est un \emph{corps de définition} de $\beta$ : pour tout
diagramme commutatif de corps
\end{itemize}

\begin{tikzcd}
L \arrow[r, "j"] \arrow[dr, dashed, "\lambda"] & M \\
k \arrow[u, "i"] \arrow[r, "i'"'] & L' \arrow[u, "j'"']
\end{tikzcd}

où $i$, $i'$ sont les homomorphismes canoniques, et tout $\beta' \in F(L')$
tel que $F(j')(\beta') = F(j)(\beta)$, on a $j'(L') \supset j(L)$ — d'où
$\lambda : L \to L'$ avec $j' \lambda = j$ — et $\beta' = F(\lambda)(\beta)$.

Ce diagramme est le (Q) de la page. La définition dit, sans supposer $F$
représentable, que $\eta$ est l'élément universel au voisinage d'un point
dont le corps résiduel est exactement $k(t)$. Elle se lit au mieux quand
$F_a$ est représentable par $X_a$ ($a$ le point de $S$ sous $t$) : $\eta$
définit $u : T \to X_a$, et, $x = u(t)$, $(T, \eta)$ est modulaire si et
seulement si $\mathcal{O}_{X_a, x} \to \mathcal{O}_{T,t}$ est un
isomorphisme (page 36). Un croquis, dans la marge de la page 34, dispose
$A \to B$, $A \to B'$, $k$, $L$, $L'$ autour d'une courbe séparant $t \in T$
de $t' \in T'$ ; plusieurs lettres y sont biffées, et on ne le reproduit pas.

\emph{N.B.} Si $\eta \in F(T)$ et $\eta' \in F(T')$ sont modulaires et qu'il
existe un diagramme (Q) avec $F(j)(\eta(t)) = F(j')(\eta'(t'))$, alors (et
seulement alors) il existe un unique $\varphi : T \to T'$ tel que $\eta =
F(\varphi)(\eta')$ : deux paires modulaires qui passent par le même point
sont isomorphes.\footnote{« et alors seulement » et « unique » sont des
lectures douteuses.}

\pagerange{36}{40}
\subsection*{La démonstration de la Proposition 3 (pages 36 à 40)}

On peut supposer $S$ noethérien. Par 1°), $M(T, \eta)$ est stable par
générisation : si $\eta$ est modulaire en $t$, l'anneau local en une
générisation de $t$ est un localisé de $\mathcal{O}_{T,t} \simeq
\mathcal{O}_{X_a,x}$. Étant constructible par 2°), il est donc ouvert. Par
(L2), pour $a \in S$ et $x \in X_a$ au-dessus de $a$, il existe un $S$-schéma
$X^{(x)}$ de type fini et un $\eta^{(x)} \in F(X^{(x)})$ modulaire en tous
les points de $X^{(x)}$.\footnote{On étale un voisinage affine de $x$ dans
$X_a$, avec son élément universel, au-dessus d'un voisinage de $a$ — c'est ce
que permet (L2) — puis on se restreint à l'ouvert de modularité. La page
écrit « il existe [\ldots] un $X^{(x)}$ noethérien, de type fini sur $S$ »,
avec deux mots illisibles.} Un premier essai pour recoller les $X^{(x)}$ —
par les ensembles $X^{(a,b)}$ des points de $X^{(a)}$ « apparentés » à un
point de $X^{(b)}$ — est encadré et barré. Il passe par un lemme :

\textbf{Lemme.} Soient $Z$ un $S$-schéma localement de type fini et $\eta \in
F(Z)$ modulaire en tout point de $Z$. Considérons $\eta : h_Z \to F$. Alors
$h_Z \times_F h_Z \subset h_{Z \times_S Z}$ est représentable par un
sous-schéma $Z'$ de $Z \times_S Z$, et les deux projections $Z'
\rightrightarrows Z$ sont des isomorphismes locaux, qui font de $Z'$ un
graphe d'équivalence dans $Z$.\footnote{La page écrit « $\eta \in F(S)$ » ;
on lit $\eta \in F(Z)$, seule lecture qui donne un sens à $h_Z \to F$. Le
lemme est suivi de l'abréviation « AQT », non résolue, et n'est pas
démontré ; la présente lecture ne le certifie pas.}

Prenons $Z = X^{(x)}$. Au-dessus de $S_a$, la relation $Z'_a
\rightrightarrows Z_a$ est la relation d'équivalence \emph{triviale} ; elle
l'est donc aussi au-dessus d'un voisinage ouvert convenable $U$ de
$a$.\footnote{Pour que la relation soit triviale au-dessus de $S_a$, il faut
que $Z_a \to X_a$ soit injectif ; c'est le cas si $X^{(x)}$ est obtenu en
étalant un ouvert de $X_a$, comme on l'a fait. C'est ce que la page appelle
prendre $X^{(x)}$ « assez petit ».} Autrement dit, quitte à prendre $X^{(x)}$
assez petit, $(X^{(x)}, \eta^{(x)})$ est « presque modulaire » : $h_{X^{(x)}}
\to F$ est une immersion ouverte. Alors « une autre proposition formelle »
permet de recoller les $X^{(x)}$ en un $(X, \xi)$ qui représente $F$, grâce à
(L2). Cette proposition n'est pas dans le dossier.

C'est sur cette phrase que finit la main de Grothendieck dans le dossier ; la
page 41, qui le clôt, est un feuillet de tapuscrit étranger. Et c'est sur cette
proposition que s'appuie, par le Lemme de la page 4, la dernière étape de la
démonstration du Théorème : les deux bouts du dossier se rejoignent là, sur
ce qui n'est pas écrit.

\end{document}
